% page 465 of the facsimile (image p-464 of the scan)
% block 167.6 x 237.7 mm ; left margin 34.4 mm
\pgc{12.700mm}{0.000mm}{
\l{\cel{53}457}
\l{}
\l{\sou{precesiono}. (\sou{astron.})\cel{2}\sou{Precesiono di la equinoxi}\cel{1}: diplaso}
\l{\cel{3}pokopa di la axo di Tero, pro quo la equinoxo printempala}
\l{\cel{3}eventas ante la ri-eso di Tero en la sama punto di lua or-}
\l{\cel{3}bito. - DEFIS.}
\l{}
\l{\sou{precipiso}. I. Pento tre eskarpa. - II. (\sou{metaf}.)Kauzo di ruino.}
\l{\cel{3}EFIS.}
\l{}
\l{\sou{precipitar}. (\sou{trans}.)\cel{2}I. Faligar bruske, de loko situita alte,}
\l{\cel{3}ad fundo. - II. (\sou{metaf.}) (1) Faligar subite, de socio-grado}
\l{\cel{3}alta ; (b) faligar subite a situeso funesta. - III. (\sou{kemio)}.}
\l{\cel{3}Retrofaligar a la fundo di la vazo qua kontenas liquido,}
\l{\cel{3}solidon qua esas dissolvita en lu. - IV. Igar irar impetuoze.}
\l{\cel{3}- DEFIS.}
\l{}
\l{\sou{precipua}. I. I. La maxim importoza, aludante (1) persono, (b)}
\l{\cel{3}kozo. - II. (\sou{gram.})\cel{1}\sou{Propoziciono precipua}\cel{1}:\cel{2}opoze a la pro-}
\l{\cel{3}poziciono incidenta. - L.}
\l{}
\l{\sou{preciza}. I. (imajo, ideo). Strikte inkluzita, limitizita, determi-}
\l{\cel{3}nita. - II. (Tranchuro, sekuro, rupturo) di qua nula saliajo}
\l{\cel{3}igas nedicidita la konturi. - DEFIRS.}
\l{}
\l{\sou{predikar}. (\sou{trans.})\cel{2}I. Anuncar lo dicita da Deo. - II. Instruk-}
\l{\cel{3}tar (ulu) per docar lo dicita da Deo. - III. Diskursar, de}
\l{\cel{3}katedro kristanismala, por la edifiko di la asistanti. -}
\l{\cel{3}DFIS.}
\l{}
\l{\sou{predikato}. (\sou{logiko}). Atributo di propoziciono. - DEFIS.}
\l{}
\l{\sou{prefaco.}\cel{2}Expozo-texto preliminara quan autoro insertas ante}
\l{\cel{3}la komenco-parto di sua verko ed en qua lu savigas olua sko-}
\l{\cel{3}po, olua karaktero, e c. - EFIS.}
\l{}
\l{\sou{prefekto.}\cel{2}I. La persono qua esas la chefo di stat-administro-}
\l{\cel{3}fako en distrikto determinita. - II. La persono qua administras}
\l{\cel{3}departmento, en Francia. - DEFIRS.}
\l{}
\l{\sou{preferar}. (\sou{trans., ulu, ulo, kam ulu, kam ulo).}\cel{2}Amar, amorar,}
\l{\cel{3}prizar, diletar plu multe (kam onu agas lo pri altru od altro).}
\l{\cel{3}- EFIS.}
\l{}
\l{\sou{prefixo}. (\sou{gram.}) Partikulo qua, lokizite avan vorto, modifikas}
\l{\cel{3}la senco di olca e formacas vorto nova. - DEFIRS.}
\l{}
\l{\sou{prefoliaciono.}\cel{2}(\sou{bot.)}\cel{2}Dispozeso-maniero di la folii en la}
\l{\cel{3}burjono. - DEFIRS.}
\l{}
\l{\sou{pregar}. (\sou{trans., ulu, pri o por}). I. Adorar la deajo, Deo,}
\l{\cel{3}demandante a lu favoro od indulgo. - II. Urjar (ulu) pri ke}
\l{\cel{3}lu grantez. - FI.}
\l{}
\l{+ prei. - Demando plu insistoza kam\cel{1}\sou{+pliz}\cel{1}- preska supliko. - E.}
}
